\chapter{Projet final : \texttt{NkAtelier}}

Vous avez maintenant traversé toute la pile : la fenêtre, les événements, le
rendu 2D, les widgets, les images, les polices, le son, la vidéo, le réseau. Ce
dernier chapitre ne présente aucune API nouvelle. Il fait une seule chose, et
c'est la plus utile : \textbf{il assemble}.

Nous allons construire, du fichier de build jusqu'à l'exécutable qui s'affiche,
une petite application appelée \texttt{NkAtelier}. Elle tient en un seul
\nkfile{main.cpp}, elle ne dépend d'\textbf{aucun fichier externe} — pas une
police, pas une image, pas un son — et elle prouve, à l'écran, que cinq modules
coopèrent.

\section{Ce que nous allons construire}

\texttt{NkAtelier} est une fenêtre avec trois panneaux.

\begin{enumerate}
  \item \textbf{Panneau « Image »} — une image \emph{fabriquée à l'exécution}
        par les primitives de dessin de NKImage, redimensionnée en vignette,
        envoyée au backend et affichée par le widget \texttt{Image}. Un curseur
        change une couleur ; l'image est refabriquée et ré-uploadée.
  \item \textbf{Panneau « Son »} — trois boutons qui jouent trois sons
        \emph{synthétisés} par \texttt{AudioGenerator}, sur le bus « UI », avec
        une case à cocher pour couper ce bus.
  \item \textbf{Panneau « Vidéo »} — au démarrage, l'application \emph{écrit}
        une petite vidéo MJPEG dans un fichier AVI, puis la \emph{relit} image
        par image et l'affiche, cadencée, dans l'interface. Boutons lecture,
        pause et retour au début.
\end{enumerate}

Le tout avec du texte rendu par une police embarquée.

\begin{nkretenir}[Pourquoi aucun actif externe]
C'est un choix pédagogique délibéré, et il vaut d'être expliqué.

Une application d'apprentissage qui exige un \texttt{logo.png} et un
\texttt{clic.wav} échoue à la première exécution, parce que le répertoire
courant n'est pas celui qu'on croit. Le lecteur passe alors une heure sur un
problème de chemin relatif au lieu d'apprendre le moteur — c'est exactement ce
que la démo du dépôt contourne avec sa liste de chemins candidats
(\nkfile{Applications/NKGuiDemo/src/NKGuiDemo/main.cpp:294-299}).

Ici, chaque ressource est \textbf{produite} :
\begin{itemize}
  \item la police vient de \texttt{NkFontEmbedded} (chapitre~6) ;
  \item l'image est dessinée par \texttt{NkImage::FillRect} et consorts
        (chapitre~5) ;
  \item le son est synthétisé par \texttt{AudioGenerator::GenerateTone}
        (chapitre~7) ;
  \item la vidéo est écrite par \texttt{NkVideoWriter} avant d'être relue
        (chapitre~8).
\end{itemize}
\textbf{L'application ne peut pas échouer faute de fichier.} C'est également ce
qui la rend facile à donner à quelqu'un d'autre.
\end{nkretenir}

\section{Étape 0 — l'arborescence}

Un module ou une application du dépôt, c'est un dossier avec un
\texttt{.jenga}. Créez :

\begin{nkcode}{Arborescence à créer}
Applications/NkAtelier/
  NkAtelier.jenga
  src/
    main.cpp
\end{nkcode}

C'est tout. Un seul fichier source ; c'est suffisant pour ce que nous faisons, et
cela évite de mélanger l'apprentissage du moteur avec celui d'une organisation
de projet.

\section{Étape 1 — le fichier de build}

\subsection{Le contenu}

Le fichier ci-dessous est modelé sur
\nkfile{Applications/NKGuiDemo/NKGuiDemo.jenga}, allégé de ce dont nous n'avons
pas besoin (la réflexion, la détection de Vulkan) et augmenté de ce qu'il nous
faut (\texttt{NKAudio}, \texttt{NKMedia}).

\begin{nkcode}{Applications/NkAtelier/NkAtelier.jenga — écrit pour le cours, modelé sur NKGuiDemo.jenga:41-95}
#!/usr/bin/env python3
# -*- coding: utf-8 -*-
"""NkAtelier — projet final du cours : image + texte + son + video dans une UI NKGui."""

from Jenga import *
from jengaconfig import *


with project("NkAtelier"):
    windowedapp()
    language("C++")
    cppdialect("C++17")
    location(".")

    files(["src/**.cpp"])

    nkentseudependson(
        ["NKGui", "NKCanvas", "NKWindow", "NKEvent", "NKGlad",
         "NKFont", "NKImage", "NKAudio", "NKMedia",
         "NKStream", "NKFileSystem", "NKLogger", "NKMath", "NKTime",
         "NKContainers", "NKMemory", "NKCore", "NKPlatform", "NKThreading"],
        extra_includes=["src", "%{NKGlad.location}/include"],
    )

    objdir("%{wks.location}/Build/Obj/%{cfg.buildcfg}-%{cfg.system}/%{prj.name}")
    targetdir("%{wks.location}/Build/Bin/%{cfg.buildcfg}-%{cfg.system}/%{prj.name}")

    # ===== Windows ============================================================
    with filter("system:Windows && !options:windows-runtime=uwp && !system:XboxSeries && !system:XboxOne"):
        windowedapp()
        usetoolchain(TC_WINDOWS)
        defines(["WIN32_LEAN_AND_MEAN", "_UNICODE", "UNICODE"])
        links([
            "user32", "gdi32", "opengl32", "dwmapi", "shell32",
            "uuid", "ole32",
            "d3d11", "d3d12", "dxgi", "d3dcompiler",
        ])

    # ===== Linux (XLib) =======================================================
    with filter("system:Linux && options:linux-backend=xlib || system:Linux && !options:linux-backend && !options:headless"):
        windowedapp()
        usetoolchain("clang-native")
        defines(["NKENTSEU_FORCE_WINDOWING_XLIB_ONLY"])
        links(["pthread", "X11", "Xext", "GL"])

    # ===== macOS ==============================================================
    with filter("system:macOS"):
        windowedapp()
        usetoolchain("clang-native")
        frameworks(["Cocoa", "QuartzCore", "OpenGL"])

    with filter("config:Debug"):
        defines(["_DEBUG", "DEBUG"])
        optimize("Off")
        symbols(True)
    with filter("config:Release"):
        defines(["NDEBUG"])
        optimize("Speed")
        symbols(False)
\end{nkcode}

\subsection{Ce qu'il faut comprendre dans ce fichier}

\begin{itemize}
  \item \texttt{windowedapp()} — application fenêtrée, pas console. Sous Windows,
        cela évite qu'une console noire s'ouvre derrière votre fenêtre ;
  \item \texttt{nkentseudependson([\dots])} — la liste des modules. Elle
        \emph{doit} contenir tout ce que vous incluez, transitivement ou non ;
  \item \texttt{extra\_includes} — \texttt{NKGlad} est nécessaire parce que le
        backend OpenGL de NkCanvas expose ses en-têtes dans les vôtres ;
  \item les blocs \texttt{with filter(\dots)} — la configuration par plateforme.
        Ne les inventez pas : recopiez-les d'une application existante qui
        fonctionne.
\end{itemize}

\begin{nkpiege}[La liste des dépendances n'est pas décorative]
Un module absent de \texttt{nkentseudependson} ne provoque pas une erreur claire
« module manquant ». Il produit soit une erreur de compilation sur un en-tête
introuvable — encore lisible —, soit une erreur d'\emph{édition de liens} sur des
symboles non résolus, qui ne dit pas quel module ajouter.

Le réflexe : quand un symbole \texttt{nkentseu::quelquechose} manque à l'édition
de liens, cherchez dans quel module il est déclaré, et ajoutez ce module.
\end{nkpiege}

\subsection{Enregistrer le projet dans l'espace de travail}

Un \texttt{.jenga} isolé n'est pas construit. Il faut l'inclure depuis le
descripteur racine, \nkfile{Nkentseu.jenga}, à côté des autres applications :

\begin{nkcode}{Nkentseu.jenga — ajout, forme identique à Nkentseu.jenga:1394-1396}
    with include("Applications/NkAtelier/NkAtelier.jenga"):

        pass
\end{nkcode}

\section{Étape 2 — le squelette : fenêtre, contexte, backend, police}

\subsection{Les inclusions}

\begin{nkcode}{Applications/NkAtelier/src/main.cpp — écrit pour le cours (modèle : NKGuiDemo/main.cpp:7-31)}
#include "NKWindow/NKWindow.h"
#include "NKWindow/NKMain.h"
#include "NKEvent/NkWindowEvent.h"
#include "NKEvent/NkMouseEvent.h"
#include "NKEvent/NkKeyboardEvent.h"

#include "NKCanvas/Core/NkContextDesc.h"
#include "NKCanvas/Core/NkGraphicsApi.h"
#include "NKCanvas/Renderer/Targets/NkRenderWindow.h"
#include "NKCanvas/UI/NkGuiCanvasBackend.h"   // pont NKGui -> NKCanvas (header-only)

#include "NKGui/NKGui.h"                      // TOUJOURS l'en-tete parapluie de NKGui
#include "NKImage/Core/NkImage.h"
#include "NKAudio/NkAudio.h"
#include "NKMedia/Video/NkVideoWriter.h"
#include "NKMedia/Video/NkVideoReader.h"

#include "NKTime/NkClock.h"
#include "NKMemory/NkUniquePtr.h"
#include "NKMemory/NKMemory.h"
#include "NKLogger/NkLog.h"
#include "NKMath/NkColor.h"

using namespace nkentseu;
using namespace nkentseu::nkgui;
using namespace nkentseu::renderer;
\end{nkcode}

Deux remarques d'ordre général.

D'abord, \textbf{on inclut l'en-tête parapluie de NKGui}
(\nkfile{NKGui/NKGui.h}) et jamais un en-tête interne — le chapitre~1 en a donné
la raison : c'est lui qui neutralise les macros de X11 sous Linux.

Ensuite, \textbf{on n'inclut pas d'en-tête parapluie pour NKMedia} : il n'en
existe pas. On prend les deux sous-en-têtes précis dont on a besoin.

\subsection{Fenêtre et cible de rendu}

\begin{nkcode}{main.cpp — écrit pour le cours (modèle : NKGuiDemo/main.cpp:240-263)}
int nkmain(const NkEntryState &state) {
    (void)state;

    NkWindow window;
    NkWindowConfig cfg;
    cfg.title = "NkAtelier - projet final du cours Nkentseu";
    cfg.width = 1180;
    cfg.height = 760;
    cfg.centered = true;
    cfg.resizable = true;
    if (!window.Create(cfg))
        return -1;

    NkContextDesc desc;
    desc.api = NkGraphicsApi::NK_GFX_API_AUTO;
    if (desc.api == NkGraphicsApi::NK_GFX_API_AUTO) {
#if defined(NKENTSEU_PLATFORM_WINDOWS)
        desc.api = NkGraphicsApi::NK_GFX_API_DX11;
#else
        desc.api = NkGraphicsApi::NK_GFX_API_OPENGL;
#endif
    }
    auto target = memory::NkMakeUnique<NkRenderWindow>(window, desc);
    if (!target || !target->IsValid())
        return -1;
\end{nkcode}

Le point d'entrée s'appelle \texttt{nkmain} et vient de
\nkfile{NKWindow/NKMain.h} : c'est lui qui gère les différences de démarrage
entre plateformes.

\begin{nkpiege}[Si rien ne s'affiche, changez de backend]
Le tableau du chapitre~1 rappelait que, à la date de sa rédaction, seul le
backend OpenGL était validé à l'exécution ; DX11 et le rasteriseur logiciel
avaient un écran noir signalé, Vulkan et DX12 des plantages
(\nkfile{Kernel/Runtime/NKCanvas/ROADMAP.md}, lignes 185-200 et 601-607).

Ces états sont datés et le dépôt bouge — mais si votre fenêtre s'ouvre noire,
\textbf{la première chose à essayer n'est pas de relire votre code} : c'est de
remplacer \texttt{NK\_GFX\_API\_DX11} par \texttt{NK\_GFX\_API\_OPENGL} et de
relancer. Cinq secondes, et cela élimine la moitié des causes possibles.
\end{nkpiege}

\subsection{Contexte NKGui et backend}

\begin{nkcode}{main.cpp — écrit pour le cours (modèle : NKGuiDemo/main.cpp:264-279)}
    auto ctxPtr = memory::NkMakeUnique<NkGuiContext>();
    if (!ctxPtr)
        return -1;
    NkGuiContext &ctx = *ctxPtr;
    ctx.Init(static_cast<int32>(cfg.width), static_cast<int32>(cfg.height));
    SetCurrentContext(&ctx);

    NkGuiCanvasBackend backend;
    if (!backend.Init(target->GetRenderer()))
        return -1;
\end{nkcode}

\textbf{L'ordre compte} : le backend reçoit le renderer, donc le renderer doit
exister. Et, souvenez-vous du chapitre~5, toutes les textures doivent être
créées \emph{après} l'initialisation du renderer — ce qui est le cas ici,
puisque nos uploads viendront plus bas.

\subsection{La police}

\begin{nkcode}{main.cpp — écrit pour le cours (modèle : NKGuiDemo/main.cpp:281-288)}
    auto fontPtr = memory::NkMakeUnique<NkGuiFont>();
    if (!fontPtr->LoadEmbedded(NkEmbeddedFontId::DroidSans, 17.f)) {
        fontPtr->LoadEmbedded(NkEmbeddedFontId::ProggyClean, 16.f);
    }
    ctx.font = fontPtr.Get();
    if (fontPtr->Valid()) {
        backend.UploadFontGray8(fontPtr->TexId(), fontPtr->pixels, fontPtr->atlasW, fontPtr->atlasH);
    } else {
        logger.Error("[NkAtelier] aucune police embarquee disponible : textes invisibles");
    }
\end{nkcode}

Trois lignes utiles et un repli, exactement comme au chapitre~6. Le \texttt{else}
n'est pas superflu : sans lui, une police manquante donnerait une interface
entièrement muette, sans le moindre indice.

\section{Étape 3 — les événements et la boucle}

\subsection{Le branchement des entrées}

\begin{nkcode}{main.cpp — écrit pour le cours (modèle : NKGuiDemo/main.cpp:326-347)}
    auto &events = NkEvents();
    bool running = true;
    events.AddEventCallback<NkWindowCloseEvent>([&](NkWindowCloseEvent *) { running = false; });
    events.AddEventCallback<NkMouseMoveEvent>([&](NkMouseMoveEvent *e) {
        ctx.input.mousePos = {static_cast<float32>(e->GetX()), static_cast<float32>(e->GetY())};
    });
    events.AddEventCallback<NkMouseButtonPressEvent>([&](NkMouseButtonPressEvent *e) {
        if (e->GetButton() == NkMouseButton::NK_MB_LEFT)
            ctx.input.mouseDown[0] = true;
    });
    events.AddEventCallback<NkMouseButtonReleaseEvent>([&](NkMouseButtonReleaseEvent *e) {
        if (e->GetButton() == NkMouseButton::NK_MB_LEFT)
            ctx.input.mouseDown[0] = false;
    });
    events.AddEventCallback<NkMouseWheelVerticalEvent>([&](NkMouseWheelVerticalEvent *e) {
        ctx.input.wheel += static_cast<float32>(e->GetDeltaY());
    });
\end{nkcode}

C'est le strict minimum pour des boutons, des curseurs et du défilement. Si vous
ajoutez un champ de saisie, il vous faudra en plus \texttt{NkTextInputEvent} et
la traduction des touches d'édition — la démo en donne le modèle complet
(\nkfile{Applications/NKGuiDemo/src/NKGuiDemo/main.cpp:367-405}).

\subsection{La boucle, et l'invariant d'ordre}

\begin{nkcode}{main.cpp — écrit pour le cours (modèle : NKGuiDemo/main.cpp:429-470)}
    NkClock clock;
    uint32 lastW = 0, lastH = 0;

    while (running && window.IsOpen()) {
        float32 dt = clock.Tick().delta;
        if (dt <= 0.f)  dt = 1.f / 60.f;
        if (dt > 0.1f)  dt = 0.1f;          // garde-fou : fenetre deplacee, machine chargee

        while (NkEvent *ev = NkEvents().PollEvent()) {
            (void)ev;                        // le travail est fait par les callbacks
        }
        if (!running)
            break;

        // Synchroniser la swapchain sur la taille de la FENETRE.
        const math::NkVec2u wsz = target->GetWindow().GetSize();
        if (wsz.x > 0 && wsz.y > 0 && (wsz.x != lastW || wsz.y != lastH)) {
            target->OnResize(wsz.x, wsz.y);
            lastW = wsz.x;
            lastH = wsz.y;
        }
        const math::NkVec2u sz = target->GetSize();
        if (sz.x > 0 && sz.y > 0) {
            ctx.viewW = static_cast<int32>(sz.x);
            ctx.viewH = static_cast<int32>(sz.y);
        }

        ctx.BeginFrame(dt);
        // ... les trois panneaux, sections suivantes ...
        ctx.EndFrame();

        target->Clear();
        backend.Submit(ctx.dl, sz.x, sz.y);
        backend.Submit(ctx.dlOverlay, sz.x, sz.y);
        target->Display();
    }
\end{nkcode}

\begin{nkretenir}[L'invariant d'ordre, une dernière fois]
\texttt{événements} $\rightarrow$ \texttt{BeginFrame(dt)} $\rightarrow$
\texttt{widgets} $\rightarrow$ \texttt{EndFrame()} $\rightarrow$ \texttt{Clear}
$\rightarrow$ \texttt{Submit} $\times 2$ $\rightarrow$ \texttt{Display}.

Les \textbf{deux} \texttt{Submit} ne sont pas une coquille : le second envoie la
couche des popups et menus, qui doit se dessiner par-dessus tout le reste.
L'oublier donne une application où les menus déroulants passent derrière les
panneaux — un bug qu'on cherche longtemps parce qu'il ressemble à un problème de
profondeur.
\end{nkretenir}

Le bornage de \texttt{dt} mérite aussi un mot. Quand l'utilisateur déplace la
fenêtre, le système bloque la boucle : le \emph{delta time} suivant peut valoir
plusieurs secondes. Toutes vos animations feraient alors un bond. Un plafond à
100~ms coûte une ligne et supprime une catégorie entière de bizarreries.

\section{Étape 4 — l'image (NKImage)}

\subsection{Fabriquer l'image}

Nous n'avons pas de fichier ; nous dessinons. Voici la fonction, écrite pour ce
cours, qui produit une image et l'envoie au backend :

\begin{nkcode}{main.cpp — écrit pour le cours (API : NkImage.h:232, 587-765, 404, 775)}
// Fabrique une image 512x512, en tire une vignette, et l'uploade sous `texId`.
// Renvoie false si quoi que ce soit echoue.
static bool ConstruireVignette(NkGuiCanvasBackend &backend, uint32 texId, float32 teinte,
                               int32 &outW, int32 &outH) {
    // 1) Une image sur la PILE : rien a liberer.
    NkImage source;
    if (!source.Create(512, 512, math::NkColor(18, 20, 26, 255), 4))
        return false;

    // 2) Dessin CPU (chapitre 5). Toutes ces primitives sont inline dans l'en-tete.
    const uint8 r = static_cast<uint8>(60.f + 195.f * teinte);
    const uint8 g = static_cast<uint8>(200.f - 120.f * teinte);
    source.FillRect(40, 40, 432, 120, math::NkColor(r, g, 90, 255));
    source.DrawRect(40, 40, 432, 120, math::NkColor(255, 255, 255, 255));
    for (int32 i = 0; i < 12; ++i)
        source.DrawCircle(256, 300, 20 + i * 12, math::NkColor(0, 220, 255, 255 - i * 18));
    source.DrawLine(0, 0, 511, 511, math::NkColor(255, 255, 255, 90));
    source.DrawLine(0, 511, 511, 0, math::NkColor(255, 255, 255, 90));

    // 3) Vignette : Resize renvoie un NkImage* -> il FAUT le liberer.
    NkImage *vignette = source.Resize(192, 192, NkResizeFilter::NK_BILINEAR);
    if (!vignette)
        return false;

    // 4) Aplatir en RGBA CONTIGU (le stride ne vaut pas forcement w*4 : chapitre 5).
    static NkVector<uint8> rgba;          // statique : on evite de reallouer a chaque appel
    const int32 w = vignette->Width();
    const int32 h = vignette->Height();
    rgba.Resize(static_cast<usize>(w) * h * 4u);
    for (int32 y = 0; y < h; ++y) {
        const uint8 *src = vignette->RowPtr(y);
        uint8 *dst = rgba.Data() + static_cast<usize>(y) * w * 4u;
        for (int32 x = 0; x < w * 4; ++x)
            dst[x] = src[x];
    }

    const bool ok = backend.UploadImageRGBA(texId, rgba.Data(), w, h);
    outW = w;
    outH = h;

    vignette->Free();                     // OBLIGATOIRE : issu d'une fabrique
    return ok;                            // ~NkImage() libere `source`
}
\end{nkcode}

Quatre points de vigilance sont concentrés dans cette fonction, et ce sont
exactement les quatre du chapitre~5 :

\begin{enumerate}
  \item \texttt{source} est sur la pile, donc rien à libérer ;
  \item \texttt{Resize} renvoie un pointeur, donc \texttt{Free()} obligatoire —
        y compris si vous ajoutez un \texttt{return} anticipé après cette ligne ;
  \item la recopie passe par \texttt{RowPtr(y)}, pas par un \texttt{memcpy}
        global ;
  \item le tampon \texttt{rgba} est \texttt{static} : on ne réalloue pas 147~Ko à
        chaque changement de curseur.
\end{enumerate}

\subsection{L'appeler, et l'afficher}

Avant la boucle :

\begin{nkcode}{main.cpp — écrit pour le cours}
    constexpr uint32 kTexVignette = 0x56494731u;   // 'VIG1'
    int32 vigW = 0, vigH = 0;
    float32 teinte = 0.35f;
    bool vignetteOk = ConstruireVignette(backend, kTexVignette, teinte, vigW, vigH);
    float32 teinteAppliquee = teinte;
\end{nkcode}

Dans la frame :

\begin{nkcode}{main.cpp — écrit pour le cours (API : NkGuiWidgets.h:39, 65, 99, 140)}
        SetNextWindowSize(ctx, 340.f, 420.f);
        if (Begin(ctx, "Image (NKImage)")) {
            Text(ctx, "Image dessinee a l'execution, sans aucun fichier.");
            Separator(ctx);

            if (vignetteOk)
                Image(ctx, kTexVignette, static_cast<float32>(vigW), static_cast<float32>(vigH));
            else
                Text(ctx, "(construction de l'image echouee)");

            Separator(ctx);
            SliderFloat(ctx, "Teinte", teinte, 0.f, 1.f);
            // On ne refabrique QUE si la valeur a change : sinon on redessinerait
            // et re-uploaderait 512x512 pixels a chaque frame pour rien.
            if (teinte != teinteAppliquee) {
                vignetteOk = ConstruireVignette(backend, kTexVignette, teinte, vigW, vigH);
                teinteAppliquee = teinte;
            }
            EndWindow(ctx);
        }
\end{nkcode}

\begin{nkretenir}[Ne refaites pas le travail à chaque image]
Le test \texttt{teinte != teinteAppliquee} est la seule chose qui sépare une
application fluide d'une application qui redessine et ré-uploade une texture
soixante fois par seconde sans raison.

C'est un réflexe à prendre : dans une interface en mode immédiat, \emph{déclarer}
un widget est bon marché, mais \emph{réagir} à un widget doit être conditionné au
changement. Les widgets qui renvoient \texttt{true} en cas de modification —
\texttt{SliderFloat}, \texttt{Checkbox}, \texttt{DragFloat} — sont faits pour
cela.
\end{nkretenir}

\section{Étape 5 — le son (NKAudio)}

\subsection{Synthétiser trois sons}

\texttt{AudioGenerator} produit des \texttt{AudioSample} sans aucun fichier.
C'est ce que font quatre applications du dépôt, dont Pong :

\begin{nkcode}{Applications/Pong/src/Pong/Audio/AudioManager.cpp:31-45}
        static AudioSample MakePaddleHit() {
            AudioSample s = AudioGenerator::GenerateTone(
                /*frequency*/ 880.0f,
                /*duration */ 0.06f,
                /*waveform */ WaveformType::SINE,
                /*sampleRate*/ 48000,
                /*amplitude*/ 0.85f);
            AdsrEnvelope env;
            env.attackTime = 0.002f; // attack immediate
            env.decayTime = 0.02f;
            env.sustainLevel = 0.5f;
            env.releaseTime = 0.03f;
            AudioGenerator::ApplyEnvelope(s, env);
            return s;
        }
\end{nkcode}

L'enveloppe ADSR n'est pas un raffinement optionnel : sans elle, un son qui
commence et s'arrête brutalement produit un clic audible — l'oreille entend la
discontinuité. Deux millisecondes d'attaque suffisent à le supprimer.

Adaptons pour notre atelier :

\begin{nkcode}{main.cpp — écrit pour le cours (API : NkAudio.h:707, 758 ; modèle : Pong/AudioManager.cpp:31-45)}
static audio::AudioSample FaireSon(float32 freq, float32 duree, audio::WaveformType forme) {
    audio::AudioSample s = audio::AudioGenerator::GenerateTone(freq, duree, forme, 48000, 0.75f);
    audio::AdsrEnvelope env;
    env.attackTime = 0.003f;
    env.decayTime = 0.02f;
    env.sustainLevel = 0.5f;
    env.releaseTime = 0.04f;
    audio::AudioGenerator::ApplyEnvelope(s, env);
    return s;
}
\end{nkcode}

\subsection{Initialiser, jouer, couper}

Avant la boucle :

\begin{nkcode}{main.cpp — écrit pour le cours (API : NkAudio.h:1091, 1104)}
    audio::AudioEngine &engine = audio::AudioEngine::Instance();
    const bool audioOk = engine.Initialize(audio::AudioEngineConfig{});   // backend AUTO
    if (!audioOk)
        logger.Error("[NkAtelier] audio indisponible : l'application continue en silence");

    audio::AudioSample sonGrave, sonMedium, sonAigu;
    if (audioOk) {
        logger.Info("[NkAtelier] device audio : {0} Hz, {1} canaux", engine.GetSampleRate(), engine.GetChannels());
        sonGrave  = FaireSon(220.f, 0.08f, audio::WaveformType::SQUARE);
        sonMedium = FaireSon(440.f, 0.07f, audio::WaveformType::SINE);
        sonAigu   = FaireSon(880.f, 0.06f, audio::WaveformType::SINE);
    }
    bool sonsUiActifs = true;
\end{nkcode}

Dans la frame :

\begin{nkcode}{main.cpp — écrit pour le cours (API : NkGuiWidgets.h:44-45 ; NkAudio.h:1127, 1266)}
        SetNextWindowSize(ctx, 340.f, 260.f);
        if (Begin(ctx, "Son (NKAudio)")) {
            if (!audioOk) {
                Text(ctx, "Moteur audio indisponible sur cette machine.");
            } else {
                Text(ctx, "Trois sons synthetises, joues sur le bus UI.");
                Separator(ctx);

                audio::VoiceParams vp;
                vp.bus = "UI";           // bus dedie : coupable sans toucher a la musique

                if (Button(ctx, "Grave (220 Hz)"))
                    engine.Play(sonGrave, vp);
                if (Button(ctx, "Medium (440 Hz)"))
                    engine.Play(sonMedium, vp);
                if (Button(ctx, "Aigu (880 Hz)"))
                    engine.Play(sonAigu, vp);

                Separator(ctx);
                if (Checkbox(ctx, "Sons d'interface", sonsUiActifs)) {
                    if (audio::NkAudioBus *bus = engine.GetBus("UI"))
                        bus->SetMute(!sonsUiActifs);
                }
            }
            EndWindow(ctx);
        }
\end{nkcode}

Remarquez que \texttt{Play} est appelé sans récupérer le \emph{handle} : ce sont
des sons courts, on ne les pilotera pas. Remarquez aussi le \texttt{if} autour de
\texttt{GetBus} : il renvoie \texttt{nullptr} si le moteur n'est pas initialisé.

\begin{nkpiege}[N'écrivez jamais \texttt{cfg.backend = DIRECTSOUND}]
Rappel du chapitre~7 : le backend DirectSound est un stub qui renvoie
\texttt{true} sans rien ouvrir
(\nkfile{Kernel/Runtime/NKAudio/src/NKAudio/NkAudioBackends.cpp:445-451}). Vous
auriez un moteur « initialisé », des voix « en lecture », et le silence complet.
Laissez \texttt{AUTO}.
\end{nkpiege}

\section{Étape 6 — la vidéo (NKMedia)}

C'est la partie la plus ambitieuse, et elle est faite d'un aller-retour : on
écrit une vidéo, puis on la relit.

\subsection{Écrire la vidéo}

\begin{nkcode}{main.cpp — écrit pour le cours (modèle : NKVideoTest/src/main.cpp:83-112)}
// Ecrit une petite video MJPEG/AVI de `nframes` images. Renvoie true si OK.
static bool EcrireVideoDemo(const char *chemin, int32 w, int32 h, int32 fps, int32 nframes) {
    media::NkVideoConfig cfg;
    cfg.width = w;
    cfg.height = h;
    cfg.fpsNum = fps;
    cfg.fpsDen = 1;
    cfg.codec = media::NkVideoCodec::MJPEG;        // le chemin le plus sur (chapitre 8)
    cfg.container = media::NkVideoContainer::AVI;
    cfg.quality = 88;

    media::NkVideoWriter vw;
    if (!vw.Open(chemin, cfg)) {
        logger.Error("[NkAtelier] impossible d'ouvrir {0} en ecriture", chemin);
        return false;
    }

    // Un SEUL tampon, reutilise. RGB24 : 3 octets par pixel, contigu, haut->bas.
    uint8 *rgb = static_cast<uint8 *>(memory::NkAlloc(static_cast<usize>(w) * h * 3));
    if (!rgb) {
        vw.Close();
        return false;
    }

    bool ok = true;
    for (int32 f = 0; f < nframes && ok; ++f) {
        // Motif : un carre qui traverse l'image + un fond qui defile.
        const int32 cx = (w - 60) * f / (nframes > 1 ? nframes - 1 : 1);
        for (int32 y = 0; y < h; ++y) {
            uint8 *ligne = rgb + static_cast<usize>(y) * w * 3;
            for (int32 x = 0; x < w; ++x) {
                const bool dansCarre = (x >= cx && x < cx + 60 && y >= h / 2 - 30 && y < h / 2 + 30);
                ligne[x * 3 + 0] = dansCarre ? 255 : static_cast<uint8>((x + f * 3) & 0xFF);
                ligne[x * 3 + 1] = dansCarre ? 140 : static_cast<uint8>((y * 2) & 0xFF);
                ligne[x * 3 + 2] = dansCarre ? 0   : 60;
            }
        }
        ok = vw.WriteFrame(rgb, media::NkVideoInputFormat::RGB24);
    }

    memory::NkFree(rgb);
    vw.Close();          // OBLIGATOIRE : ecrit l'index AVI. Sans lui : fichier illisible.
    return ok;
}
\end{nkcode}

Notez que \texttt{Close()} est appelé \textbf{sur tous les chemins}, y compris
quand \texttt{ok} est devenu \texttt{false}. C'est la règle du chapitre~8 : sans
finalisation, le fichier existe et n'est lisible par personne.

Notez aussi la construction de la ligne : \texttt{rgb + y * w * 3}. Ici, pas de
\emph{stride} aligné — nous fabriquons nous-mêmes un tampon contigu, c'est ce que
\texttt{WriteFrame} attend. Le \emph{stride} est une notion de \texttt{NkImage},
pas de \texttt{NkVideoWriter}.

\subsection{Relire la vidéo}

Avant la boucle :

\begin{nkcode}{main.cpp — écrit pour le cours (API : NkVideoReader.h:54-66)}
    constexpr uint32 kTexVideo = 0x56494431u;      // 'VID1'
    const char *kCheminVideo = "nkatelier_demo.avi";

    media::NkVideoReader lecteur;
    bool videoOk = false;
    double videoFps = 25.0;
    int32 videoW = 0, videoH = 0, videoNb = -1;

    if (EcrireVideoDemo(kCheminVideo, 320, 240, 25, 75)) {   // 3 secondes
        if (lecteur.Open(kCheminVideo)) {
            const media::NkVideoReaderInfo &in = lecteur.Info();
            videoW = in.width;
            videoH = in.height;
            videoNb = in.frameCount;
            videoFps = (in.fps > 1.0) ? in.fps : 25.0;
            videoOk = true;
            logger.Info("[NkAtelier] video : {0} {1} {2}x{3} @ {4} fps",
                        in.container.CStr(), in.codec.CStr(), in.width, in.height, in.fps);
        } else {
            logger.Error("[NkAtelier] video ecrite mais illisible : {0}", kCheminVideo);
        }
    }

    media::NkVideoFrame trame;      // DECLAREE HORS BOUCLE : son buffer est reutilise
    float32 accumulateur = 0.f;
    bool videoEnLecture = true;
    bool videoFinie = false;
    int32 videoIndex = -1;
\end{nkcode}

Le fait que la lecture soit tentée juste après l'écriture est en soi un test :
si \texttt{Open} échoue, c'est que l'écriture a mal tourné.

\subsection{Cadencer et afficher}

\begin{nkcode}{main.cpp — écrit pour le cours (modèle : NKCode/Shell/NkVideoViewer.h:153-170)}
        SetNextWindowSize(ctx, 400.f, 420.f);
        if (Begin(ctx, "Video (NKMedia)")) {
            if (!videoOk) {
                Text(ctx, "Video indisponible.");
            } else {
                // ── Cadencement ────────────────────────────────────────────
                const float32 dureeImage = 1.f / static_cast<float32>(videoFps);
                if (videoEnLecture && !videoFinie) {
                    accumulateur += dt;
                    if (accumulateur > dureeImage * 4.f)
                        accumulateur = 0.f;          // evite le rattrapage explosif
                    while (accumulateur >= dureeImage) {
                        accumulateur -= dureeImage;
                        if (lecteur.ReadFrame(trame)) {
                            videoIndex = trame.index;
                            // Une seule image decodee -> un seul upload.
                            backend.UploadImageRGBA(kTexVideo, trame.rgba.Data(), trame.width, trame.height);
                        } else {
                            videoFinie = true;       // fin du fichier
                            break;
                        }
                    }
                }

                // ── Affichage ──────────────────────────────────────────────
                if (videoIndex >= 0)
                    Image(ctx, kTexVideo, static_cast<float32>(videoW), static_cast<float32>(videoH));
                else
                    Text(ctx, "(premiere image en attente)");

                Separator(ctx);
                Text(ctx, videoFinie ? "Etat : terminee" : (videoEnLecture ? "Etat : lecture" : "Etat : pause"));

                BeginHBox(ctx);
                if (Button(ctx, videoEnLecture ? "Pause" : "Lecture"))
                    videoEnLecture = !videoEnLecture;
                if (Button(ctx, "Debut")) {
                    if (lecteur.SeekFrame(0)) {
                        videoFinie = false;
                        accumulateur = 0.f;
                        videoIndex = -1;
                    }
                }
                EndHBox(ctx);
            }
            EndWindow(ctx);
        }
\end{nkcode}

\begin{nkretenir}[Les trois règles de la lecture vidéo, réunies]
\begin{enumerate}
  \item \textbf{Un seul upload par image affichée.} La boucle \texttt{while}
        peut décoder plusieurs trames si l'application a pris du retard, mais on
        n'uploade que celle qu'on va montrer ;
  \item \textbf{l'accumulateur est borné} (\texttt{> dureeImage * 4}) : sans ce
        garde-fou, un décrochage de deux secondes déclencherait le décodage de
        cinquante images d'un coup, ce qui aggraverait le décrochage ;
  \item \textbf{\texttt{NkVideoFrame} est déclarée hors boucle.} Son
        \texttt{NkVector} garde sa capacité ; la déclarer à l'intérieur
        provoquerait une allocation par image.
\end{enumerate}
\end{nkretenir}

\begin{nkpiege}[Le rebouclage n'est exact que sur les formats non compressés en temps]
\texttt{SeekFrame(0)} fonctionne bien ici parce que notre vidéo est en MJPEG :
chaque image est indépendante. Sur du H.264, \texttt{SeekFrame(n)} ne fait que
revenir à l'image autonome précédente, et il faut redécoder vers l'avant
jusqu'à la cible (chapitre~8). Notre choix du MJPEG n'était donc pas seulement
un choix de prudence : il simplifie aussi la navigation.
\end{nkpiege}

\begin{nkpiege}[L'upload répété dépend de \texttt{NkTexture::Update}]
Regardez l'implémentation du pont : \texttt{UploadImageRGBA} crée la texture au
premier appel, puis appelle \texttt{tex->Update(...)} aux suivants
(\nkfile{Kernel/Runtime/NKCanvas/src/NKCanvas/UI/NkGuiCanvasBackend.h:94-117}).
C'est exactement le bon comportement — on ne recrée pas la texture par image.

Mais souvenez-vous de l'avertissement de Pong (chapitre~6) : « \texttt{Update}
exige \texttt{gTextureBackend.Update} qui n'est pas câblé sur tous les backends
$\rightarrow$ échec silencieux » (\nkfile{Applications/Pong/src/Pong/Render/FontAtlas.cpp:107-109}).

Si votre vignette statique s'affiche mais que la vidéo reste figée sur sa
première image, la cause est probablement là — et la solution est de changer de
backend graphique, pas de réécrire votre code de lecture.
\end{nkpiege}

\section{Étape 7 — la fermeture}

C'est le moment où l'on récolte ce qu'on a semé dans les cinq chapitres
précédents. Chaque module a son ordre, et ils s'imbriquent.

\begin{nkcode}{main.cpp — écrit pour le cours (invariants : NkAudio.h:1114 ; NkAudioPlayer/main.cpp:258-259)}
    // ── Fermeture, dans cet ordre exact ───────────────────────────────────────

    // 1) Video : fermer le lecteur (il possede son etat interne).
    lecteur.Close();

    // 2) Audio : ARRETER LE MOTEUR AVANT de liberer les samples.
    if (audioOk) {
        engine.Shutdown();
        audio::AudioLoader::Free(sonGrave);
        audio::AudioLoader::Free(sonMedium);
        audio::AudioLoader::Free(sonAigu);
    }

    // 3) Interface.
    SetCurrentContext(nullptr);
    ctx.Shutdown();
    return 0;
    // Les NkUniquePtr (target, ctxPtr, fontPtr) se detruisent ici, dans l'ordre
    // inverse de leur declaration. Le backend et l'atlas de police suivent.
}
\end{nkcode}

\begin{center}
\begin{tabular}{@{}lll@{}}
\toprule
\textbf{Module} & \textbf{Ce qu'on ferme} & \textbf{Piège si l'ordre est inversé} \\
\midrule
NKMedia & \texttt{lecteur.Close()} & fichier verrouillé \\
NKAudio & \texttt{Shutdown()} puis \texttt{Free()} & lecture après libération \\
NKImage & rien (pile) ou \texttt{Free()} & fuite ou double libération \\
NKFont  & rien (l'atlas possède tout) & double libération si fusion \\
NKGui   & \texttt{ctx.Shutdown()} & — \\
\bottomrule
\end{tabular}
\end{center}

\begin{nkretenir}[Le principe unique derrière tous ces ordres]
\textbf{On éteint d'abord ce qui \emph{tourne}, ensuite ce qui est \emph{lu}.}

Le moteur audio a un fil qui lit vos échantillons ; le gestionnaire de connexion
a un fil qui lit vos tampons ; l'enregistreur vidéo a un fil qui lit vos images.
Dans les trois cas, la fonction d'arrêt joint le fil — et ce n'est qu'après
qu'on a le droit de libérer ce qu'il lisait.

Si vous ne deviez retenir qu'une phrase de ces six chapitres, ce serait
celle-là : elle explique l'ordre de fermeture de NKAudio, de NKNetwork et de
NKMedia d'un seul coup.
\end{nkretenir}

\section{Compiler et lancer}

\begin{nkcode}{Ligne de commande (voir README.md:147-163)}
jenga build --target NkAtelier --config Debug
\end{nkcode}

L'exécutable atterrit dans
\nkfile{Build/Bin/Debug-Windows/NkAtelier/} — ou
\texttt{Debug-Linux}, \texttt{Debug-macOS} selon la plateforme, suivant le
\texttt{targetdir} que nous avons écrit dans le \texttt{.jenga}.

Pour la version optimisée :

\begin{nkcode}{Ligne de commande}
jenga build --target NkAtelier --config Release
\end{nkcode}

\begin{nkpiege}[Le répertoire courant au lancement]
Notre application écrit \texttt{nkatelier\_demo.avi} \textbf{dans le répertoire
courant}, pas à côté de l'exécutable. Selon que vous la lancez depuis un
terminal, depuis un explorateur de fichiers ou depuis un environnement de
développement, ce répertoire diffère.

Ce n'est pas grave ici — nous écrivons \emph{et} lisons au même endroit dans la
même exécution — mais c'est exactement le mécanisme qui fait qu'un
\texttt{assets/logo.png} « introuvable » l'est seulement dans certaines
conditions de lancement. Si vous voulez voir le fichier produit, lancez
l'application depuis un terminal, dans un dossier que vous avez choisi.

Souvenez-vous aussi que les auto-tests de NKMedia écrivent, eux aussi, dans le
répertoire courant (chapitre~8).
\end{nkpiege}

\section{Le catalogue des silences}

Voici la partie la plus utile de ce chapitre. Toutes les pannes ci-dessous ont
un point commun : \textbf{aucune ne produit de message d'erreur}. Gardez cette
liste sous la main.

\begin{center}
\begin{tabular}{@{}p{5.4cm}p{8.4cm}@{}}
\toprule
\textbf{Symptôme} & \textbf{Cause à vérifier en premier} \\
\midrule
Fenêtre noire, rien du tout &
Backend graphique. Essayez \texttt{NK\_GFX\_API\_OPENGL}. \\
Tout s'affiche sauf les textes &
Atlas de police non uploadé, ou \texttt{ctx.font} non affecté. \\
Les menus passent \emph{derrière} &
Second \texttt{Submit(ctx.dlOverlay, \dots)} oublié. \\
L'image apparaît « penchée » &
\texttt{memcpy} global au lieu de \texttt{RowPtr(y)}. \\
L'image est vide &
Texture créée avant \texttt{renderer.Initialize()}. \\
Aucun son, mais tout indique que ça joue &
Backend \texttt{DIRECTSOUND} demandé explicitement. \\
Grésillement ou plantage à la fermeture &
\texttt{Free(sample)} appelé avant \texttt{engine.Shutdown()}. \\
Le fichier vidéo est illisible &
\texttt{Close()} non appelé sur le \texttt{NkVideoWriter}. \\
La vidéo reste sur la première image &
\texttt{NkTexture::Update} non câblé dans ce backend. \\
La vidéo est à l'envers &
\texttt{flipVertical} avec un framebuffer OpenGL. \\
Vignette de mauvaise qualité malgré Lanczos &
\texttt{NK\_LANCZOS3} retombe sur du bilinéaire. \\
Aucun pair ne se connecte &
Poignée de main asynchrone : attendre \texttt{ConnectedPeerCount()}. \\
Erreurs d'édition de liens \texttt{\_\_imp\_socket} &
\texttt{ws2\_32} absent du \texttt{links()}. \\
\bottomrule
\end{tabular}
\end{center}

\begin{nkretenir}[La méthode générale]
Ces silences ont tous la même parade, et c'est une habitude à prendre plutôt
qu'une astuce : \textbf{journalisez ce que vous chargez, et journalisez aussi ce
que vous ne chargez pas}.

Toutes les fonctions d'initialisation de ce chapitre ont un \texttt{else} qui
écrit une ligne. Cela paraît excessif quand tout fonctionne. Le jour où quelque
chose manque, cette ligne vous fait gagner une heure — et c'est précisément ce
que fait la démo du dépôt : « \texttt{Image demo introuvable (cwd) -> section
image vide} »
(\nkfile{Applications/NKGuiDemo/src/NKGuiDemo/main.cpp:422}).
\end{nkretenir}

\section{Le programme, vu de haut}

Récapitulons la structure complète, dans l'ordre du fichier :

\begin{nkcode}{Applications/NkAtelier/src/main.cpp — plan d'ensemble}
// --- inclusions (etape 2) ---

// --- fonctions utilitaires ---
static bool ConstruireVignette(backend, texId, teinte, &w, &h);   // etape 4
static audio::AudioSample FaireSon(freq, duree, forme);           // etape 5
static bool EcrireVideoDemo(chemin, w, h, fps, nframes);          // etape 6

int nkmain(const NkEntryState &state) {
    // 1. fenetre + cible de rendu                                   etape 2
    // 2. contexte NKGui + backend                                   etape 2
    // 3. police embarquee + upload atlas                            etape 2
    // 4. callbacks d'evenements                                     etape 3
    // 5. image : premiere construction + upload                     etape 4
    // 6. audio : Initialize + synthese des trois sons               etape 5
    // 7. video : ecriture du fichier puis ouverture en lecture      etape 6

    while (running && window.IsOpen()) {
        // dt + pompage des evenements + synchro swapchain           etape 3
        ctx.BeginFrame(dt);
        //   panneau Image                                           etape 4
        //   panneau Son                                             etape 5
        //   panneau Video (cadencement + upload + affichage)        etape 6
        ctx.EndFrame();
        // Clear + Submit x2 + Display                               etape 3
    }

    // fermeture : video, puis audio (Shutdown avant Free), puis UI  etape 7
}
\end{nkcode}

Environ 300 lignes en tout. C'est peu pour cinq modules — et c'est le point.
\textbf{La difficulté du moteur n'est pas dans la quantité de code à écrire, elle
est dans les ordres à respecter} : ordre d'initialisation, ordre de la frame,
ordre de fermeture.

\section{Pour aller plus loin}

Trois extensions, par difficulté croissante. Chacune réutilise un chapitre
précédent sans rien introduire de nouveau.

\subsection{Enregistrer ce que l'atelier affiche}

Ajoutez un bouton « Enregistrer » qui, à chaque frame,
capture la fenêtre et pousse l'image dans un \texttt{NkVideoRecorder} en mode
MJPEG :

\begin{nkcode}{Esquisse — écrit pour le cours (modèle : NKViewportDemo/main.cpp:329-362)}
if (enregistrement && target->CaptureToImage(capture) && capture.IsValid()) {
    if (!rec.IsRecording())
        rec.Begin("nkatelier_capture.mp4", capture.Width(), capture.Height(),
                  30, 1, 22, 32, media::NkRecorderCodec::MJPEG, 92);
    if (rec.IsRecording())
        rec.PushVideo(capture.Pixels(), media::NkVideoInputFormat::RGBA32);
}
// ... et, sur TOUS les chemins de sortie :
if (rec.IsRecording())
    rec.End();
\end{nkcode}

Attention aux deux pièges du chapitre~8 : en mode MJPEG l'audio est ignorée
silencieusement, et \texttt{End()} doit être appelé même si la fenêtre est fermée
brutalement.

\subsection{Ajouter le réseau}

Ajoutez à \texttt{NkAtelier} un quatrième panneau : deux boutons « Héberger » et
« Rejoindre », et une liste de messages. Chaque clic sur un bouton de son envoie
un message aux autres instances, qui jouent le même son.

Vous aurez besoin d'ajouter \texttt{NKNetwork} aux dépendances du
\texttt{.jenga} \emph{et} \texttt{ws2\_32} au \texttt{links()} sous Windows. Le
chapitre~9 donne la structure à suivre : une petite classe de session, un
\texttt{Tick(dt)} par frame, et une interface qui ne voit qu'un état.

\subsection{Faire du décodage un travail de fond}

Notre vidéo fait 320 par 240 pixels : le décodage est instantané. Essayez avec
une vidéo de 1920 par 1080, et vous verrez la boucle hoqueter.

La solution — décoder sur un fil, uploader sur le fil principal — est celle du
chapitre~5, appliquée à la vidéo. C'est un excellent exercice, mais commencez par
\textbf{mesurer} : entourez \texttt{ReadFrame} d'un chronomètre mural et
regardez. Et souvenez-vous de l'avertissement du chapitre~8 : mesurez le temps
mur, pas le \emph{PTS} de la vidéo.

\begin{nkbilan}
Une application complète n'est pas la somme des chapitres : c'est leur
\emph{ordre}. Créer avant d'employer, fermer dans le sens inverse de
l'ouverture, et n'ouvrir une ressource qu'après ce dont elle dépend. Le
« catalogue des silences » est la partie la plus utile de ce chapitre : dans
chaque module, il existe des façons d'échouer qui ne disent rien --- un fichier
absent, un identifiant nul, une passe sans attachement --- et la seule parade est
de \emph{vérifier ce qu'on reçoit} plutôt que de supposer que l'absence
d'erreur vaut succès.
\end{nkbilan}

\section{Exercices}

\begin{nkdemo}[1 — Faire tomber chaque panneau, un par un]
Pour chacun des quatre modules, cassez volontairement une chose et notez ce que
vous observez :
\begin{itemize}
  \item n'appelez pas \texttt{UploadFontGray8} ;
  \item supprimez le \texttt{vignette->Free()} et lancez l'application quelques
        minutes en bougeant le curseur de teinte (surveillez la mémoire) ;
  \item inversez \texttt{Shutdown()} et \texttt{Free()} à la fermeture ;
  \item supprimez le \texttt{vw.Close()} dans \texttt{EcrireVideoDemo}.
\end{itemize}
Écrivez, pour chaque cas, la phrase que vous auriez aimé lire dans un journal.
Puis ajoutez-la vraiment à votre code. Vous venez d'écrire votre propre
catalogue des silences.
\end{nkdemo}

\begin{nkpratique}[2 — Une vraie vignette de fichier]
Ajoutez un cinquième panneau qui charge une image \emph{depuis le disque} avec
\texttt{NkImage::Load}, en essayant plusieurs chemins candidats comme le fait la
démo du dépôt. Affichez le format d'origine (\texttt{SourceFormat()}), les
dimensions, le nombre de canaux et le \emph{stride}.

Faites-en une vignette de 128 pixels de large en préservant le rapport d'aspect,
et affichez côte à côte l'original réduit et la vignette. Puis ajoutez un bouton
qui sauvegarde la vignette en PNG et un autre qui tente de la sauvegarder en
WebP. Affichez les deux valeurs de retour côte à côte. Que constatez-vous ?
\end{nkpratique}

\begin{nkpratique}[3 — Le mélangeur]
Ajoutez trois \texttt{SliderFloat} — Master, UI, Music — reliés aux bus
correspondants par \texttt{GetBus(\dots)->SetVolume(\dots)}, en n'appelant le
\emph{setter} que lorsque la valeur change.

Ajoutez ensuite un son long en boucle sur le bus « Music »
(\texttt{vp.looping = true}) et vérifiez, à l'oreille, la formule du
chapitre~7 : le volume effectif est le produit du volume de la voix, de celui du
bus, et de tous ses parents jusqu'au Master. Vérifiez enfin que couper « UI » ne
touche pas à « Music ».
\end{nkpratique}

\begin{nkpratique}[4 — Une chaîne complète, de l'image à la vidéo et retour]
Modifiez \texttt{EcrireVideoDemo} pour qu'elle ne dessine plus ses pixels à la
main, mais qu'elle utilise une \texttt{NkImage} et ses primitives de dessin —
comme \texttt{ConstruireVignette} — puis qu'elle envoie
\texttt{image.Pixels()} à \texttt{WriteFrame}.

Attention : \texttt{WriteFrame} attend un tampon \textbf{contigu}, et
\texttt{NkImage::Pixels()} a un \emph{stride} aligné sur 4 octets. Choisissez
donc des dimensions telles que \texttt{stride == width * bpp}, ou bien recopiez
ligne par ligne. Vérifiez votre raisonnement en affichant les deux valeurs avant
d'écrire la première image.

Relisez ensuite la vidéo produite et affichez-la : vous aurez fait circuler les
mêmes pixels à travers NKImage, le codec JPEG, le conteneur AVI, le décodeur, et
enfin le GPU. C'est toute la chaîne de ce cours en une seule fonction.
\end{nkpratique}

\begin{nkpratique}[5 — Le vôtre]
Reprenez \texttt{NkAtelier} et transformez-le en quelque chose qui vous est
utile : une visionneuse de dossier d'images, un métronome, un petit lecteur de
séquences, un banc d'essai de vos propres filtres d'image.

Une seule contrainte, et c'est celle de tout ce cours : \textbf{pour chaque
fonction du moteur que vous emploierez, vérifiez dans la source qu'elle fait
bien ce que son nom annonce}. Vous savez maintenant où regarder — le
\texttt{.cpp}, pas le commentaire ; le \texttt{ROADMAP.md}, pas le
\texttt{.jenga} — et vous savez que la réponse est parfois « non ».
\end{nkpratique}

\begin{nkquestions}
\begin{enumerate}
\item Dans quel ordre crée-t-on fenêtre, contexte, backend et police ? Pourquoi
      pas un autre ?
\item Pourquoi ferme-t-on dans l'ordre inverse de l'ouverture ?
\item Citez trois « silences » du catalogue, et ce qui les rend difficiles à
      voir.
\item Que faut-il faire quand une taille de fenêtre change ?
\item Qu'est-ce qui, dans ce programme, ne pourrait pas être écrit sans les
      chapitres précédents ?
\end{enumerate}
\end{nkquestions}

\begin{nklogique}
Votre application se lance, la fenêtre s'ouvre, tout est noir. Aucun message
d'erreur.

\begin{enumerate}
\item Énumérez au moins cinq causes possibles, en les rangeant par étage de la
      pile.
\item Pour chacune, dites ce qui la distinguerait des autres --- une observation
      dont le résultat \emph{diffère} selon la cause.
\item Dans quel ordre feriez-vous ces vérifications ? Justifiez par le
      \emph{coût} de chacune, pas par sa probabilité --- et expliquez pourquoi ce
      critère est le bon quand on ne sait rien.
\end{enumerate}
\end{nklogique}
